\documentclass{beamer}

\usepackage{slmacros}

\usepackage{hyperref}
\usepackage{tikz}

\usetheme[secheader]{Boadilla}

\slideFrontPage{Introduction to Distributed Data Systems}{}

%\date{\today}

\begin{document}

\maketitle


%
% Command to show the table of content at the beginning of each section
%

\AtBeginSection{
\begin{frame}<beamer>
  \frametitle{Outline}
  {\tableofcontents[subsectionstyle=show/show/hide,sectionstyle=show/shaded]}
\end{frame}
}

\section{Introduction}


\begin{frame}
\frametitle{Outline}

\vfill

\alert{Guidelines and principles for distributed data management}

\begin{itemize}
  \item Basics of distributed systems (networks, performance, principles).
    \item Failure management and distributed transactions.  
     \item Properties of a distributed system
     \item Specifics of peer-to-peer networks
     \item A powerful and resilient distribution method: consistent hashing
     \item Case study: GFS, a distributed file system.
\end{itemize}

\end{frame}

\section{Overview of distributed data management principles}


\begin{frame}
\frametitle{Distributed systems}

A \alert{distributed system} is an application that coordinates the actions of
  several computers to achieve a specific task. 
  
  \vfill
  This coordination is achieved by exchanging \alert{messages}
  which are  pieces of data that convey
   some information. 
   
   $\Rightarrow$ ``shared-nothing'' architecture -> no shared memory, no shared
   disk.
   
   \vfill 
   
   The system relies on a network that connects the
   computers and handles the routing of messages.
   
   $\Rightarrow$ Local area networks (LAN), Peer to peer (P2P) networks, \ldots
    
    \vfill
    
    \alert{Client} (nodes) and \alert{Server} (nodes) are communicating
    \alert{software} components: we assimilate them with the machines they run
    on.
\end{frame}

\begin{frame}
\frametitle{LAN-based infrastructure: clusters of machines}

Three communication levels: ``racks'', clusters, and groups of clusters. 

\figSlideWithSize{network}{9cm}

\end{frame}


\begin{frame}
\frametitle{Example: data centers}

Typical setting of a Google data center.

\begin{enumerate}
  \item $\approx$ 40 servers per rack;
  \item $\approx$ 150 racks per data center (cluster);

  \item $\approx$ 6,000 servers per data center;
    \item how many clusters? Google's secret, and constantly evolving \ldots
\end{enumerate}

Rough estimate: 150-200 data centers? 1,000,000 servers? 

\end{frame}


\begin{frame}
\frametitle{P2P infrastructure: Internet-based communication}

 Nodes, or ``peers'' communicate with messages
  sent over the Internet network.
  
  \vfill
  
\figSlideWithSize{internet}{9cm}

\vfill

 The physical route may consist of 10 or more forwarding
  messages,  or ``hops''.
  
  \vfill
  
  Suggestion: use the \texttt{traceroute} utility to 
  check the route between your laptop and a Web site of your choice.
\end{frame}

\begin{frame}
\frametitle{Performance}

\begin{center}
\begin{small}

\begin{tabular}{l|p{5cm}|p{5cm}}
{\bf Type} & {\bf Latency} & \textbf{Bandwidth}\\
\hline
Disk & $\approx 5 \times
10^{-3}\mathrm{s}$ (5
millisec.); & At best 100 MB/s\\ 
LAN &  $\approx 1-2 \times
10^{-3}\mathrm{s}$ (1-2
millisec.);  &  $\approx 1 \mathrm{GB/s}$ (single rack);  $\approx 100
\mathrm{MB/s}$ (switched); \\
Internet &  Highly variable. Typ. 10-100 ms.;  &  Highly variable.
Typ. a few MB/s.;\\
\hline
\end{tabular}

Bottom line (1): it is approx. one order of magnitude faster to exchange main
memory data between 2 machines in a data center, that to read on the disk.

Bottom line (2): exchanging through the Internet is slow and unreliable with
respect to LANs.

\end{small}
\end{center}
\end{frame}


\begin{frame}
\frametitle{Distribution, why?}

% A two-columns presentation
\begin{tabular}{l|p{6cm}}
   \parbox{5cm}{

  \textbf{Sequential access.}  It takes 166 minutes (more than 2 hours and a
  half) to read a 1 TB disk. 
  
  \vfill
  \textbf{Parallel access.} With 100 disks, assuming that the disks work in
  parallel and sequentially:  about 1mn 30s.
 \vfill
 
 \textbf{Distributed access.} With 100 computers, each disposing of its own
 local disk: each CPU processes its own dataset.  
   }  % Inclusion of an XML file
  & 
  \parbox{6cm}{\figSlideWithSize{bigpic-distrstorage}{5.5cm} }
\end{tabular}
\vfill

% A Beamer box
\begin{beamerboxesrounded}{Scalability}
The latter solution is \emph{scalable}, by adding new computing resources.
\end{beamerboxesrounded}

\end{frame}


\begin{frame}
\frametitle{What you should remember: performance of data-centric distr.
systems}

   \begin{enumerate}
      \item disk transfer rate is a bottleneck for large scale data management;
      parallelization and distribution of the data on many machines is a means
      to eliminate this bottleneck;
      \item \emph{write once, read many}: a distributed storage system is
      appropriate for large files that are written once and then
      repeatedly scanned;
       \item \emph{data locality}: bandwidth is a scarce resource, and program
      should be ``pushed'' near the data they must access to.
     \end{enumerate}

A distr. system also gives an opportunity to reinforce the security
of data with \alert{replication}.
      
\end{frame}

\subsection{Data replication and consistency}

\begin{frame}
\frametitle{What is it about?}

\alert{Replication}: a mechanism that copies data item located on a machine
$A$ to a remote machine $B$

\vfill

$\Rightarrow$ one obtains \alert{replica}

\vfill

\alert{Consistency}: ability of a system to behave as if the transaction of each
      user always run in isolation from other transactions, and never fails.

\vfill
 Example: shopping basket in an e-commerce application.
\vfill

$\Rightarrow$ difficult in centralized systems because of multi-users and
concurrency.

\vfill

$\Rightarrow$ even more difficult in distributed  systems because of 
replica.

\end{frame}

\begin{frame}
\frametitle{Some illustrative scenarios}

\figSlideWithSize{replication}{11cm}

\end{frame}



\begin{frame}
\frametitle{Consistency management in distr. systems}

\alert{Consistency}: essentially, ensures that the system faithfully reflects
the actions of a user.
\begin{itemize}
  \item \alert{Strong consistency} (ACID properties) --  requires a
  (slow) synchronous replication, and possibly heavy locking mechanisms.
  \item \alert{Weak consistency} -- accept to serve some requests with outdated
  data.
  \item \alert{Eventual consistency} -- same as before, but the system is
  guaranteed to converge towards a consistent state based on the last version.
\end{itemize}

\vfill

In a system that is not eventually consistent, \alert{conflicts} occur
and the application must take care of \alert{data reconciliation}: given the two
conflicting copies, determine the new current one.

\vfill

Standard RDBMS favor consistency over availability -- one of the reasons (?) of
the 'NoSQL' trend.

\end{frame}


\begin{frame}
\frametitle{Achieving strong consistency -- The 2PC protocol}

2PC = Two Phase Commit. The algorithm of choice to ensure ACID properties
in a distributed setting.

\vfill

\alert{Problem}: the update operations may occur on distinct
 servers $\{S_1, \ldots, S_n\}$, called \emph{participants}.

\end{frame}

\begin{frame}
\frametitle{Two Phase Commit: Overview}

 \figSlideWithSize{2pc}{11cm}


\end{frame}


\begin{frame}
\frametitle{Exercises and questions}

You are a customer using an e-commerce application which is known to be
eventually consistent (e.g., Amazon \ldots):
\begin{enumerate}
  \item Show a scenario where you buy an item, but this item does not appear
  in your basket.
  \item You reload the page: the item appears. What happened?
  \item You delete an item from your command, and add another one: the basket
  shows both items. What happened?
  \item Will the situation change if you reload the page?
  \item Would you expect both items to disapear in an e-commerce application?
\end{enumerate}

\end{frame}


\section{Properties of a distributed system}

\begin{frame}
\frametitle{Properties of a distributed system: (i) scalability}

\alert{Scalability} refers to the ability of a system to 
 continuously evolve in order to support an evergrowing amount
 of tasks. 

\figSlideWithSize{scalability}{11cm}

A scalable system should (i) distribute evenly the task load to all
participants, and (ii) ensure a negligible  distribution management cost.
\end{frame}


\begin{frame}
\frametitle{Properties of a distributed system: (ii) efficiency}

Two usual measures of its efficiency are the \alert{response time} (or latency)
 which denotes the delay to obtain the first item, and the \alert{throughput}
 (or bandwidth)
 which denotes the number of items delivered in a given period unit (e.g., a
 second).

\vfill

Unit costs:

 \begin{enumerate}
   \item  \emph{number of messages} 
   globally sent by the nodes of the system, regardless of the message size;
   \item  \emph{size of messages} representing  the volume of data
   exchanges.
 \end{enumerate}
 
\end{frame}

\begin{frame}
\frametitle{Properties of a distributed system: (iii) availability}

\alert{Availability} is the capacity of of a system to limit as much as possible 
its latency.
Involves several aspects:
\begin{itemize}
  \item \alert{Failure detection.}
  \begin{itemize}
    \item  monitor the participating nodes to detect
  failures as early as possible (usually via ``heartbeat'' messages); 
  \item design quick restart protocols.
\end{itemize}

  \item \alert{Replication} on several nodes.
  
  Replication may be \alert{synchronous} or \alert{asynchronous}. 
  In the later case, a Client \textit{write()} returns before the operation
  has been completed on each site.

\end{itemize} 

\end{frame}

\section{Failure management}

\begin{frame}
\frametitle{Failure recovery in centralized DBMSs}

Common principles:

\begin{enumerate}
  \item The \alert{state} of a (data) system is the set of item 
  committed by the application.
  \item Updating ``in place'' is considered as inefficient because of disk
  seeks.
  \item Instead, update are written in main memory \emph{and} in a sequential
  \alert{log file}.
  \item Failure? The main memory is lost, but all the committed  transactions
  are in the log: a \textsc{Redo} operations is carried out when the system
  restarts.
\end{enumerate}

$\Rightarrow$ implemented in all DBMSs.
\end{frame}

\begin{frame}
\frametitle{Failure and distribution}

First: when do we kown that a component is in failed state?

\vfill
$\Rightarrow$ periodically send message to each participant.

\vfill
Second: does the centralized recovery still hold?

\vfill

Yes, providing the log file is accessible \ldots

\end{frame}


\begin{frame}
\frametitle{Distributed logging}


 \figSlideWithSize{recovery-basics}{11cm}

Note: distributed logging can be asynchronous (efficient, risky)
or asynchronous (just the opposite).

\end{frame}

\section{Consistent hashing}


\begin{frame}
\frametitle{Basics: Centralized Hash files}

The collection consists of $(key, value)$ pairs. A \alert{hash function} evenly
distributes the values in \alert{buckets} w.r.t. the key.

\vfill 

\figSlideWithSize{hash-file}{9cm}

\vfill

This is the basic, \alert{static}, scheme: the number of buckets is fixed.

\vfill

\alert{Dynamic hashing} extends the number of buckets as the collection grows --
the most popular method is \alert{linear hashing}.

\end{frame}

\begin{frame}
\frametitle{Issues with hash structures distribution}

Straighforward idea: everybody uses the same hash function, and buckets are
replaced by servers. 

\vfill

Two issues:

\begin{itemize}
  \item \alert{Dynamicity}. At Web scale, we must be able to add or remove
  servers at any moment.
\item \alert{Inconsistencies}. It is very hard to ensure that all participants
share an accurante view of the system (e.g., the hash function).
\end{itemize}

\vfill

Some solutions:
\begin{itemize}
  \item \alert{Distributed linear hashing}: sophisticated scheme that allows
  Client nodes to use an outdated image of the has file; guarantees eventual
  convergence.
  \item \alert{Consistent hashing}: to be presented next.
  
  NB: consistent hashing is used in several systems, including Dynamo
  (Amazon)/Voldemort (Open Source), and P2P structures, e.g. Chord.
\end{itemize}
\end{frame}


\begin{frame}
\frametitle{Consistent hashing}

Let $N$ be the number of servers. The following functions
$$hash(key) \to modulo (key, N) = i$$
maps a pair $(key, value)$ to server $S_i$.

\vfill

\alert{Fact}: if $N$ changes, or if a client uses an invalid value for $N$, the
mapping becomes inconsistent.

\vfill

With \alert{Consistent hashing}, addition or removal of an instance does
not significantly change the mapping of keys to servers.

\begin{itemize}
  \item a simple, non-mutable hash function $h$ maps \alert{both} the keys ad
   the servers IPs to a large address space $A$ (e.g., $[0, 2^{64}-1$]);
  \item $A$ is organized as a ring, scanned in clockwise order;
  \item if $S$ and $S'$ are two adjacent servers on the ring:  
  all the keys in range $]h(S), h(S')]$ are mapped to $S'$.
\end{itemize}
\end{frame}


\begin{frame}
\frametitle{Illustration}

Example: item $A$ is mapped to server IP1-2; item B to server \ldots

\vfill

\figSlideWithSize{consist-hashing}{7cm}

\vfill

A server is added or removed? A \alert{local} re-hashing is sufficient.
\end{frame}


\begin{frame}
\frametitle{Some (really useful) refinements}

What if a server fails?  How can we balance the load? 

\vfill

% A two-columns presentation
\begin{tabular}{l|p{5cm}}
   \parbox{6.5cm}{


\alert{Failure} $\Rightarrow$ use replication; put a copy on the next machine
(on the ring), then on the next after the next, and so on.
  
  \vfill
  
\alert{Load balancing} $\Rightarrow$ map a server to several points on the ring
(virtual nodes)

\begin{itemize}
  \item the more points, the more load received by a server;
  \item also useful if the server fails: data rellocation is more evenly
  distributed.
  \item also useful in case of heterogeneity (the rule in large-scale systems).
\end{itemize}
  
   }  % Inclusion of an XML file
  & 
  \parbox{5cm}{\figSlideWithSize{ch-refinement}{5cm} }
\end{tabular}

\end{frame}



\begin{frame}
\frametitle{Distributed indexing based on consistent hashing}

%Consistent hashing has been initialized designed for caching system. 
%\vfill

Main question: \alert{where is the hash
directory (servers locations)}? Several possible answers:


\begin{itemize}
  \item  \alert{On a specific (``Master") node}, acting as a load
  balancer. Example: caching systems.
  
  $\Rightarrow$ raises scalability issues.
 
  \item \alert{Each node  records its successor on the ring}.
  
    $\Rightarrow$ may require $O(N)$ messages for routing queries -- not
    resilient to failures.
 
  \item  \alert{Each node  records $\log N$ carefully chosen other nodes.}
 
    $\Rightarrow$ ensures $O(\log N)$ messages for routing queries -- convenient
    trade-off for highly dynamic networks (e.g., P2P)

 \item \alert{Full duplication of
  the hash directory at each node.} 
  
   $\Rightarrow$ ensures $1$ message for routing  -- heavy maintenance protocol
   which can be achieved through \alert{gossiping} (broadcast of any event
   affecting the network topology).

\end{itemize}

\end{frame}


\begin{frame}
\frametitle{Case study: Dynamo (Amazon)}

A distributed system that targets high availability (your shopping cart is
stored there!).

\begin{itemize}
  \item Duplicates and maintains the hash directory at \alert{each
   node} via \alert{gossiping} --
  queries can be routed to the correct server with 1 message.
  
  \item The hosting server replicates $N$ (application parameter) copies of its
  objects on the $N$ \alert{distinct} nodes that follow $S$ on the ring.
  
  \item Propagates updates \alert{asynchronously} $\to$ may result in 
     update conflicts, solved by the application at read-time.
     
     \item Use a fully distributed failure detection mechanism (failure are
     detected by individual nodes when then fail to communicate with others)
\end{itemize}

An Open-source version is available at
\textit{http://project-voldemort.com/}
\end{frame}


\section{Case study: large scale distributed file systems (GFS)}


\begin{frame}
\frametitle{History and development of GFS}

Google File System, a paper published in 2003 by Google Labs at OSDI.

\vfill

Explains the design and architecture of a distributed system apt at 
serving very large data files; internally used by Google for storing documents
collected from the Web.

\vfill

Open Source versions have been developed at once: Hadoop File System (HDFS),
and Kosmos File System (KFS).

\end{frame}


\begin{frame}
\frametitle{The problem}

Why do we need a distributed file system in the first place?

\vfill

Fact: standard NFS (left part) does not meet  scalability requirements
(what if \file{file1} gets really big?).

\figSlideWithSize{dfs-bigpic}{9cm}

Right part: GFS/HDFS storage, based on (i) a virtual file namespace, and
(ii)  partitioning of files in ``chunks''.

\end{frame}

\begin{frame}
\frametitle{Architecture}

A \alert{Master node} performs administrative tasks, while \alert{servers}
store ``chunks'' and send them to Client nodes.

\figSlideWithSize{gfs}{8cm}

The Client maintains a \alert{cache} with chunks locations, and directly
communicates with servers.

\end{frame}

\begin{frame}
\frametitle{Technical details}

\begin{itemize}
  \item The architecture works best for very large files (e.g., several
  Gigabytes), divided in large (64-128 MBs) chunks.
 
     $\Rightarrow$ this limits the metadata information served by the Master.
     
  \item Each server implements recovery and replication techniques (default: 3
  replicas).
  
  \item (\alert{Availability}) The Master sends heartbeat messages to servers,
  and initiates a replacement when a failure occurs.
  
  \item (\alert{Scalability}) The Master is a potential single point of failure;
  its protection relies on distributed recovery techniques for all changes that affect the
  file namespace.
\end{itemize}

\end{frame}


\begin{frame}
\frametitle{Workflow of a \textit{write()} operation (simplified)}

The following figure shows a non-concurrent \textit{append()} operation.

\figSlideWithSize{gfs-write}{8.5cm}

In case of concurrent appends to a chunk, the primary replica 
assigns serial numbers to the mutation, and coordinates the secondary replicas.

\end{frame}

\begin{frame}
\frametitle{Namespace updates: distributed recovery protocol}

Extension of standard techniques for recovery (left: centralized; right:
distributed).

\figSlideWithSize{recovery-basics}{9cm}

If a node fails, the replicated log file can be used to recover the last
transactions on one of its mirrors.

\end{frame}

\end{document}

http://www-poleia.lip6.fr/~doucet/CoursBDRA/BDRA3-4-P2P.pdf








